% ==============================================================================
% reports/chapters/07_ket_luan.tex
% CHƯƠNG 7: KẾT LUẬN VÀ HƯỚNG PHÁT TRIỂN
% ==============================================================================

\chapter{Kết Luận và Hướng Phát Triển}
\label{chap:ket_luan}

Chương này tổng kết toàn bộ các kết quả đạt được từ quá trình nghiên cứu thực nghiệm, trả lời trực tiếp bốn câu hỏi nghiên cứu đã đặt ra, đồng thời phân tích các mối đe dọa đến tính hợp lệ, đề xuất các hướng mở rộng trong tương lai và báo cáo chi tiết bảng phân công nhiệm vụ của các thành viên trong nhóm.

\section{Tổng kết kết quả và trả lời các câu hỏi nghiên cứu}
\label{sec:rq_answers}

Thực nghiệm trên CICIDS2017 với 12 tổ hợp mô hình và kịch bản dữ liệu cho phép trả lời bốn câu hỏi nghiên cứu:

\begin{enumerate}
    \item \textbf{Câu hỏi nghiên cứu 1 (RQ1 -- Đánh giá trên subtype SSH vắng mặt trong tập fit theo thời gian):}\\
    Khi mô hình học máy dạng cây chỉ được huấn luyện trên tập Train chứa biến thể \FTP, khả năng phát hiện biến thể \SSH\ xuất hiện ở giai đoạn thời gian sau suy giảm nghiêm trọng. Trong kịch bản thực nghiệm chính có cổng mạng (\texttt{time/with\_port}), Random Forest chỉ đạt $\FOne = \RFTestFOne$ với $\text{SSH Recall} = \RFTestSSHRecall$ (chỉ nhận diện chính xác \RFTestTP\ trên tổng số \DataTestSSH\ luồng tấn công SSH), trong khi Decision Tree và XGBoost hoàn toàn không phát hiện được bất kỳ luồng SSH nào ($\FOne = \DTTestFOne$, $\text{SSH Recall} = \DTTestSSHRecall$). Mặc dù tập kiểm định có chứa \DataValidationSSH\ mẫu SSH tham gia vào quá trình lựa chọn siêu tham số, việc mô hình không được huấn luyện trên mẫu SSH nào trong tập Train vẫn khiến tỷ lệ phát hiện biến thể mới trên tập kiểm thử ở mức rất thấp.

    \item \textbf{Câu hỏi nghiên cứu 2 (RQ2 -- So sánh các giao thức phân tách):}\\
    Random Split và Time-based Split cho kết quả khác biệt. Trên tập kiểm thử, F1 của Random Forest lần lượt là $\RFRandomFOne$ và $\RFTestFOne$; tỷ lệ phát hiện SSH lần lượt là \RFRandomSSHRecall\ và \RFTestSSHRecall. Đây là so sánh giữa hai giao thức có thành phần subtype khác nhau: Train/Test của Random Split chứa SSH, Test còn có FTP; Time-based Train không có SSH và Test chỉ có SSH. Vì vậy, không thể quy toàn bộ chênh lệch cho cách chia hoặc cho tương quan giữa các flow. Sự phụ thuộc giữa các flow cùng chiến dịch có thể góp phần làm ước lượng Random Split lạc quan, nhưng nghiên cứu không xác lập trực tiếp rò rỉ nhãn hay đặc trưng.

    \item \textbf{Câu hỏi nghiên cứu 3 (RQ3 -- Vai trò của cổng dịch vụ và các đặc trưng hành vi):}\\
    \texttt{Destination Port} đóng vai trò là một lối tắt phân loại chi phối mạnh mẽ quyết định của mô hình. Phân tích SHAP chỉ ra cổng dịch vụ có giá trị đóng góp trung bình tuyệt đối $\text{Mean } |\text{SHAP}| = \SHAPTopOneVal$, gấp \SHAPRatioTopOneTwo\ lần đặc trưng xếp thứ hai. Khi loại bỏ cổng dịch vụ khỏi không gian huấn luyện (\texttt{without\_port}), mô hình nhận diện được thêm một lượng nhỏ các luồng SSH dựa trên hình thái thống kê gói tin (tỷ lệ SSH Recall tăng lên mức $\RFTestSSHRecall\ \text{--} \RFTestWithoutPortSSHRecall$), với F1 cao nhất đạt $\RFTestWithoutPortFOne$ ở Random Forest (14/\DataTestSSH\ luồng SSH). Tuy nhiên, hiệu năng tổng thể vẫn ở mức rất thấp, chứng minh rằng sự thất bại trong việc phát hiện SSH không thể được giải thích chỉ bằng riêng yếu tố cổng dịch vụ. Sự khác biệt về cơ chế hoạt động và phân phối đặc trưng hình thái gói tin giữa FTP và SSH, cùng với yếu tố thời gian, phù hợp với hiện tượng dịch chuyển phân phối quan sát được; mức đóng góp riêng của từng yếu tố cần được kiểm tra thêm.

    \item \textbf{Câu hỏi nghiên cứu 4 (RQ4 -- So sánh hai pipeline XGBoost):}\\
    Lần chạy tuning độc lập truyền Train + Validation vào quy trình tìm kiếm; theo thiết kế của scikit-learn \cite{scikit_learn_refit}, \texttt{refit=True} fit mô hình tốt nhất trên toàn bộ dữ liệu đầu vào này, gồm \DataValidationSSH\ mẫu SSH ở Validation. Benchmark chuẩn và pipeline tuning cũng khác nhau về số lượng/loại đặc trưng, tiền xử lý và siêu tham số, nên chênh lệch Test F1 từ $\XGBTestFOne$ đến $\XGBRefitTestFOne$ không cô lập tác động của refit và không được diễn giải như một ước lượng nhân quả. Kết quả minh họa tầm quan trọng của việc công bố rõ ranh giới fit và các khác biệt giữa pipeline.
\end{enumerate}

\section{Đánh giá tính hợp lệ và các giới hạn của đề tài}
\label{sec:threats_validity}

Bảng \ref{tab:threats_validity} tổng hợp các mối đe dọa đến tính hợp lệ của nghiên cứu và các biện pháp kiểm soát khoa học tương ứng đã được nhóm áp dụng xuyên suốt đề tài.

% ==============================================================================
% reports/tables/manual/tab_threats_validity.tex
% ==============================================================================
\begin{table}[htbp]
\centering
\caption{Tổng hợp các mối đe dọa đến tính hợp lệ của nghiên cứu và giải pháp kiểm soát.}
\label{tab:threats_validity}
\adjustbox{max width=\textwidth}{%
\begin{tabular}{l>{\raggedright\arraybackslash}p{5.5cm}>{\raggedright\arraybackslash}p{6.5cm}}
\toprule
\textbf{Khía cạnh hợp lệ} & \textbf{Nguy cơ tiềm ẩn} & \textbf{Biện pháp kiểm soát trong nghiên cứu} \\
\midrule
\textbf{Tính hợp lệ bên trong}\newline(\textit{Internal Validity}) & Sai lệch do tái dựng timestamp từ CSV, nhãn SSH trong Validation tham gia chọn cấu hình và phụ thuộc giữa các quan sát qua ranh giới tập. & Nêu rõ giờ CSV 1--5 được ánh xạ thành 13--17 và mốc cắt chọn sau khi khảo sát nhãn/thời gian; không có PCAP gốc để xác minh. Nhãn SSH trong Validation tham gia chọn cấu hình nhưng không fit trọng số benchmark. Purge/Embargo giảm nguy cơ flow vượt ranh giới; SHA-256 chỉ xác minh byte CSV đầu vào. \\
\midrule
\textbf{Tính hợp lệ xây dựng}\newline(\textit{Construct Validity}) & Nhầm lẫn giữa phân loại luồng đã hoàn tất và hệ thống NIDS thời gian thực. & Tuyên bố minh bạch phạm vi phân loại luồng đã hoàn tất (Completed-Flow Classification); không đo lường độ trễ phát hiện gói tin; không tuyên bố triển khai trực tuyến. \\
\midrule
\textbf{Tính hợp lệ bên ngoài}\newline(\textit{External Validity}) & Dữ liệu mô phỏng trong môi trường phòng thí nghiệm (CICIDS2017). & Phân tích cơ chế suy giảm tổng quát hóa; nhấn mạnh yêu cầu đánh giá trên lưu lượng mạng đa dạng và hạ tầng thực tế. \\
\midrule
\textbf{Tính hợp lệ thống kê}\newline(\textit{Statistical Validity}) & Tỷ lệ tấn công thấp ($2\% \text{--} 5\%$) gây nhiễu chỉ số Accuracy. & Sử dụng F1-Score, đường cong Precision-Recall, AP và FPR làm hệ quy chiếu cốt lõi; loại bỏ sự phụ thuộc vào nghịch lý độ chính xác (Accuracy Paradox). \\
\bottomrule
\end{tabular}%
}
\end{table}


Bên cạnh các kết quả đạt được, nhóm nghiên cứu cũng thẳng thắn ghi nhận các giới hạn phạm vi của bài tập lớn:
\begin{itemize}
    \item \textbf{Phạm vi dữ liệu và thời gian}: Toàn bộ thực nghiệm sử dụng CICIDS2017 Tuesday. Timestamp được tái dựng hồi cứu từ giờ trong CSV, ánh xạ giờ 1--5 thành 13--17; không có đồng hồ PCAP gốc để xác minh từng timestamp, và các mốc cắt được chọn sau khi khảo sát nhãn/thời gian. SHA-256 xác minh byte của CSV đầu vào, không xác nhận tính chính xác của timestamp tái dựng. Purge/Embargo giảm nguy cơ flow vượt ranh giới nhưng không xác minh đồng hồ nguồn hay chứng minh độc lập hoàn toàn. Các kết luận cần được kiểm chứng thêm trước khi khái quát hóa cho môi trường mạng khác.
    \item \textbf{Cơ chế phân loại luồng đã đóng}: Hệ thống hoạt động theo cơ chế phân loại luồng mạng đã hoàn tất dựa trên các đặc trưng thống kê tổng hợp của công cụ CICFlowMeter \cite{sharafaldin2018generating}. Hệ thống không thực hiện đo lường độ trễ xử lý gói tin ở mức thời gian thực và không thay thế trực tiếp cho các hệ thống ngăn ngừa xâm nhập mạng (IPS) chặn gói tin trực tuyến như Snort \cite{roesch1999snort} hay Suricata \cite{suricata_oisf}.
    \item \textbf{Tính tái lập}: Toàn bộ mã nguồn, cấu hình siêu tham số và mã băm SHA-256 đã được khóa và lưu vết độc lập trong tệp \texttt{frozen.json}, đảm bảo khả năng tái lập kết quả trong cùng điều kiện môi trường thực nghiệm.
\end{itemize}

\section{Hướng phát triển và mở rộng}
\label{sec:future_work}

Từ các phát hiện thực nghiệm và giới hạn đã xác định, nhóm đề xuất một số hướng phát triển tiếp theo cho đề tài:
\begin{enumerate}
    \item \textbf{Mở rộng không gian tấn công Brute Force}: Đánh giá năng lực tổng quát hóa của mô hình trên các giao thức dò quét mật khẩu phổ biến khác như RDP (Cổng 3389), SMB (Cổng 445), Telnet (Cổng 23) và xác thực HTTP POST.
    \item \textbf{Nghiên cứu cơ chế phân loại luồng sớm}: Thay vì đợi luồng mạng kết thúc hoàn toàn (vốn có thể kéo dài hàng chục giây), nghiên cứu trích xuất đặc trưng thống kê chỉ từ $N$ gói tin đầu tiên ($N \in \{5, 10, 20\}$) của mỗi phiên kết nối nhằm giảm thiểu tối đa độ trễ phát hiện trong thực tế.
    \item \textbf{Mô hình kết hợp (Hybrid Architecture)}: Kết hợp giữa hệ thống phát hiện xâm nhập dựa trên luật và chữ ký tĩnh để xử lý các thuộc tính cổng/giao thức tĩnh, kết hợp với các mô hình học máy phi giám sát hoặc phương pháp phát hiện bất thường để nhận diện các hành vi bất thường trong phân phối thời gian và độ dài gói tin.
    \item \textbf{Thực nghiệm trên lưu lượng mạng thực tế}: Thu thập và kiểm thử mô hình trên dữ liệu lưu lượng mạng thực tế tại trung tâm dữ liệu hoặc mạng nội bộ trường đại học để đánh giá mức độ thích ứng trước các biến động lưu lượng thực tế.
\end{enumerate}

\vspace{\baselineskip}
\section{Bảng phân công nhiệm vụ và đóng góp của thành viên}
\label{sec:member_contributions}

Để đảm bảo tính minh bạch, công bằng và hiệu quả trong quá trình thực hiện bài tập lớn, nhóm sinh viên (Lớp \textbf{68CS2}) đã phân chia nhiệm vụ chuyên trách theo từng mảng công việc cụ thể. Bảng phân công nhiệm vụ chi tiết kèm sản phẩm bàn giao và cam kết mức độ hoàn thành của từng thành viên nhóm được trình bày trang trọng tại phần đầu báo cáo (xem Bảng Phân công nhiệm vụ, trang \pageref{page:contributions}). Bảng \ref{tab:member_contributions} dưới đây tóm lược lại các mảng trách nhiệm chuyên trách chính và tỷ lệ đóng góp của từng thành viên trong nhóm:

\begin{table}[!h]
\centering
\caption{Bảng phân công nhiệm vụ và tỷ lệ đóng góp của các thành viên trong nhóm.}
\label{tab:member_contributions}
\adjustbox{max width=\textwidth}{%
\begin{tabular}{c l c >{\raggedright\arraybackslash}p{6.8cm} c}
\toprule
\textbf{STT} & \textbf{Họ và tên} & \textbf{Mã số SV} & \textbf{Nhiệm vụ chuyên trách} & \textbf{Đóng góp} \\
\midrule
1 & Nguyễn Việt Tuấn Anh & 0204068 & Nghiên cứu cơ sở lý thuyết NIDS; khảo sát giao thức FTP/SSH và tấn công Patator \cite{lorenzo_patator}; xây dựng pipeline tiền xử lý dữ liệu và làm sạch ranh giới Purge/Embargo. & 20\% \\
\midrule
2 & Trần Phùng Đức Anh & 0204168 & Thiết kế quy trình phân tách dữ liệu theo thời gian và ngẫu nhiên; huấn luyện, tinh chỉnh siêu tham số mô hình Decision Tree và Random Forest; tổng hợp báo cáo và đồ thị thực nghiệm. & 20\% \\
\midrule
3 & Nguyễn Vũ Dũng & 0205968 & Cấu hình và thực thi mô hình XGBoost; triển khai thực nghiệm bóc tách cổng mạng; phân tích thực nghiệm trường hợp điển hình về thiết kế pipeline và cờ \texttt{refit=True}. & 20\% \\
\midrule
4 & Nguyễn Thu Hương & 0209168 & Triển khai phân tích khả năng giải thích bằng SHAP TreeExplainer \cite{lundberg2017unified}; trích xuất và đối chiếu cấu trúc rẽ nhánh Decision Tree; xây dựng bảng biểu thống kê đặc trưng. & 20\% \\
\midrule
5 & Võ Nguyễn Anh Kiệt & 0209968 & Xây dựng hệ thống kiểm thử tự động; quản lý và kiểm soát toàn vẹn chuỗi băm SHA-256 (\texttt{frozen.json}); rà soát các mối đe dọa đến tính hợp lệ và biên tập tài liệu kỹ thuật LaTeX. & 20\% \\
\bottomrule
\end{tabular}%
}
\end{table}
